\subsection{MATRICES OVER A RING}

The most important matrices in Mathematics are matrices over a ring A. The set \(A^{I \times K}\) of matrices over A corresponding to indexing sets I, K then has canonically an (A, A)-bimodule structure (§ 1, no. 14).

For every ordered pair \((i,k)\in\mathbf{I}\times\mathbf{K}\) , let \(E_{ik}\) be the matrix \((a_{jl})\) such that \(a_{ik}=1\) and \(a_{jl}=0\) for \((j,l)\neq(i,k)\) ; the \(E_{ik}\) are called the matrix units in the set of matrices \(\mathbf{A}^{\mathbf{I}\times\mathbf{K}}\) ; if I and K are finite, they form the canonical basis of this set for its left or right A-module structure (§ 1, no. 11). Clearly

\[
{ } ^ { t } E _ { i k } = E _ { k i } .
\]

Unless otherwise mentioned, the product \(M'M''\) of two matrices over A (assumed to be defined) will always be understood to be relative to the multiplication \((x,y)\mapsto xy\) in A (or, as is also said, will be ``calculated in A''). Then we have (no. 2) the associativity and distributivity formulae

\[
(X Y) Z = X (Y Z)\tag{7}
\]

\[
\left\{ \begin{array}{l} X (Y + Z) = X Y + X Z \\ (X + Y) Z = X Z + Y Z \end{array} \right.\tag{8}
\]

for three matrices X, Y, Z over A, whenever the sums and products appearing in these formulae are defined.

In particular, if \(E_{ik}\) (resp. \(E_{kl}'\) , \(E_{il}''\) ) are the matrix units in \(\mathbf{A}^{\mathbf{I} \times \mathbf{K}}\) (resp. \(\mathbf{A}^{\mathbf{K} \times \mathbf{L}}\) , \(\mathbf{A}^{\mathbf{I} \times \mathbf{L}}\) ) respectively, with \(\mathbf{I} = [1, p]\) , \(\mathbf{K} = [1, q]\) , \(\mathbf{L} = [1, r]\) , we obtain the formulae

\[
\left\{ \begin{array}{l l} E _ {i k} E _ {j l} ^ {\prime} = 0 & \text { if } k \neq j \\ E _ {i k} E _ {k l} ^ {\prime} = E _ {i l} ^ {\prime \prime}. \end{array} \right.\tag{9}
\]

Let \(\mathbf{A}^0\) be the opposite ring of \(\mathbf{A}\) and let \(a * b (= ba)\) denote the product of \(a\) and \(b\) in \(\mathbf{A}^0\) ; then, for two matrices \(X, Y\) over \(\mathbf{A}\) whose product is defined,

\[
{ } ^ { t } ( X Y ) = { } ^ { t } Y * { } ^ { t } X\tag{10}
\]

where on the right hand side \(^{t}Y\) and \(^{t}X\) are considered as matrices with elements in \(A^{0}\) ; when A is commutative, then

\[
{ } ^ { t } ( X Y ) = { } ^ { t } Y . { } ^ { t } X\tag{11}
\]

PROPOSITION 1. Let A, B be two rings and \(M = (m_{ik})_{(i,k) \in I \times K}\) and

\[
M ^ {\prime} = \left(m _ {i k} ^ {\prime}\right) _ {(i, k) \in \mathbf {I} \times \mathbf {K}}
\]

two matrices with finite indexing sets over an (A, B)-bimodule G. Suppose that for every matrix unit \(L = (a_{i})_{i \in I}\) with one row and elements in A and every matrix unit \(C = (b_{k})_{k \in K}\) with one column and elements in B, L.M.C = L.M'.C (the products being calculated via the external laws of the (A, B)-module G); then \(M = M'\) .

If \(L\) is taken to be the matrix unit \((a_s)\) with \(a_i = 1\) , \(a_s = 0\) for \(s \neq i\) , and \(C\) the matrix unit \((b_t)\) with \(b_k = 1\) , \(b_t = 0\) for \(t \neq k\) , the products \(L.M.C\) and \(L.M'.C\) are matrices with a single element respectively equal to \(m_{ik}\) and \(m'_{ik}\) .

Let A, B be two rings and \(\sigma: A \to B\) a homomorphism.

For every matrix \(M = (m_{\iota\kappa})\) over \(\mathbf{A}\) , we shall denote by \(\sigma(M)\) the matrix \((\sigma(m_{\iota\kappa}))\) over \(\mathbf{B}\) ; clearly \(\sigma(aM) = \sigma(a)\sigma(M)\) , \(\sigma(Ma) = \sigma(M)\sigma(a)\) for \(a \in \mathbf{A}\) , also \(\sigma(^t M) = ^t(\sigma(M))\) and

\[
\left\{ \begin{array}{c} \sigma (M + M ^ {\prime}) = \sigma (M) + \sigma (M ^ {\prime}) \\ \sigma (M M ^ {\prime}) = \sigma (M) \sigma (M ^ {\prime}) \end{array} \right.\tag{12}
\]

when the operations considered are defined, the products on the left and right hand sides of (12) being calculated in A and B respectively. When \(\sigma\) is denoted by \(x \mapsto x^{\sigma}\) , we write \(M^{\sigma}\) instead of \(\sigma(M)\) .

Consider in particular an anti-endomorphism \(\sigma\) of A, that is a homomorphism of A to the opposite ring \(\mathbf{A}^0\) , or a mapping of A into itself such that

\[
\sigma (a + a ^ {\prime}) = \sigma (a) + \sigma (a ^ {\prime}), \quad \sigma (a a ^ {\prime}) = \sigma (a ^ {\prime}) \sigma (a)
\]

for all \(a, a'\) in \(\mathbf{A}\) ; then, for two matrices \(M, M'\) over \(\mathbf{A}\) whose product \(MM'\) is defined,

\[
\sigma (M M ^ {\prime}) = ^ {t} (\sigma (^ {t} M ^ {\prime}). \sigma (^ {t} M))\tag{13}
\]

where the products on the two sides are calculated in A; this follows immediately from (10) and (12).

